\documentclass[draft]{agujournal2025}
\usepackage{amsmath,amssymb}
\usepackage{booktabs,longtable,array,tabularx}
\usepackage{url,xurl}
\usepackage{hyperref}
\providecommand{\tightlist}{\setlength{\itemsep}{0pt}\setlength{\parskip}{0pt}}
\setlength{\emergencystretch}{3em}
% Still required under agujournal2025.cls: removing \sloppy reintroduces four
% overfull \hbox warnings in the body text.
\sloppy
\DeclareUnicodeCharacter{00A7}{\S{}}
\DeclareUnicodeCharacter{00B0}{\ensuremath{^\circ}}
\DeclareUnicodeCharacter{00B1}{\ensuremath{\pm}}
\DeclareUnicodeCharacter{00B2}{\textsuperscript{2}}
\DeclareUnicodeCharacter{00B3}{\textsuperscript{3}}
\DeclareUnicodeCharacter{00B7}{\ensuremath{\cdot}}
\DeclareUnicodeCharacter{00D7}{\ensuremath{\times}}
\DeclareUnicodeCharacter{00FC}{\"u}
\DeclareUnicodeCharacter{0177}{\^y}
\DeclareUnicodeCharacter{0394}{\ensuremath{\Delta}}
\DeclareUnicodeCharacter{03A3}{\ensuremath{\Sigma}}
\DeclareUnicodeCharacter{03B2}{\ensuremath{\beta}}
\DeclareUnicodeCharacter{03B3}{\ensuremath{\gamma}}
\DeclareUnicodeCharacter{03B4}{\ensuremath{\delta}}
\DeclareUnicodeCharacter{03BA}{\ensuremath{\kappa}}
\DeclareUnicodeCharacter{03BB}{\ensuremath{\lambda}}
\DeclareUnicodeCharacter{03C0}{\ensuremath{\pi}}
\DeclareUnicodeCharacter{03C1}{\ensuremath{\rho}}
\DeclareUnicodeCharacter{03C3}{\ensuremath{\sigma}}
\DeclareUnicodeCharacter{03C4}{\ensuremath{\tau}}
\DeclareUnicodeCharacter{03C6}{\ensuremath{\varphi}}
\DeclareUnicodeCharacter{2013}{--}
\DeclareUnicodeCharacter{2014}{---}
\DeclareUnicodeCharacter{2016}{\ensuremath{\|}}
\DeclareUnicodeCharacter{2026}{\ldots{}}
\DeclareUnicodeCharacter{2074}{\textsuperscript{4}}
\DeclareUnicodeCharacter{207B}{\textsuperscript{-}}
\DeclareUnicodeCharacter{2081}{\ensuremath{_1}}
\DeclareUnicodeCharacter{2113}{\ensuremath{\ell}}
\DeclareUnicodeCharacter{2192}{\ensuremath{\rightarrow}}
\DeclareUnicodeCharacter{2208}{\ensuremath{\in}}
\DeclareUnicodeCharacter{2212}{\ensuremath{-}}
\DeclareUnicodeCharacter{221A}{\ensuremath{\surd}}
\DeclareUnicodeCharacter{2248}{\ensuremath{\approx}}
\DeclareUnicodeCharacter{2264}{\ensuremath{\leq}}
\DeclareUnicodeCharacter{2265}{\ensuremath{\geq}}
\DeclareUnicodeCharacter{27E8}{\ensuremath{\langle}}
\DeclareUnicodeCharacter{27E9}{\ensuremath{\rangle}}

\begin{document}

% agujournal2025.cls requires \journalname after \begin{document}; the 2019
% class took it in the preamble.  In the submission (non-``published'') branch it
% feeds the "manuscript submitted to ..." running head.
\journalname{Water Resources Research}

\title{Reported skill in daily river water-temperature prediction shrinks under strong baselines and whole-region holdout}

% AUTHOR BLOCK TO BE COMPLETED.  Names, affiliations, ORCIDs, the author count,
% and the corresponding author are deliberately unfilled and are not invented by
% this generator.  Replace with one \authors entry carrying the final agreed
% order in the 2025 form --- ``Name\affil{1}\thanks{funding}, Name\affil{2}''
% --- one \affiliation per distinct affiliation, and a verified institutional
% address.  The signed intake schema is
% docs/FAIR_SUBMISSION_READINESS_AND_TEMPLATES.md section 2.
\authors{[AUTHOR LIST TO BE COMPLETED]\affil{1}}
\affiliation{1}{[AFFILIATION 1 TO BE COMPLETED --- department or laboratory, institution, city, postcode, country]}
\affiliation{2}{[AFFILIATION 2 TO BE COMPLETED --- add or delete affiliation lines to match the final author list]}

% \authoraddr is the 2025 corresponding-author interface and replaces the 2019
% \correspondingauthor{name}{email} pair.  Both are emitted on purpose: in the
% ``published'' branch \authoraddr is typeset in the first-page margin and
% \correspondingauthor does not exist, while in the submission branch used here
% \authoraddr is defined as a no-op and only \correspondingauthor reaches the
% page.  Keep the two in agreement when the placeholders are filled.
\authoraddr{[CORRESPONDING AUTHOR TO BE COMPLETED --- full name, department,
institution, street, city, state, postcode, country
([INSTITUTIONAL E-MAIL TO BE COMPLETED])]}
\correspondingauthor{[CORRESPONDING AUTHOR TO BE COMPLETED]}{[INSTITUTIONAL E-MAIL TO BE COMPLETED]}

% Present because agujournaltemplate.tex carries them.  The class defines both as
% no-ops in either branch; AGU sets the running heads itself.
\authorrunninghead{}
\titlerunninghead{}

% Three-argument macro in agujournal2025.cls, not the 2019 ``keypoints''
% environment.  All three slots are always supplied; ``{}'' would mark an
% unused one, but the canonical Markdown is required to carry exactly three.
\keypoints%
    {Station-median skill at a seven-day lead is +0.251 against persistence but +0.038 against damped persistence at 120 U.S. gauges.}
    {A gradient-boosted tree with site identity has the lowest station-median RMSE at 1, 3, and 7 days on identical prediction keys.}
    {Holding out whole hydrologic regions instead of random sites lowers three-day skill against persistence from +0.187 to +0.116.}

% Required by agujournaltemplate.tex after the front matter.  It is a no-op in
% the submission branch and typesets the Wiley title page under ``published''.
\maketitle

\begin{abstract}
Daily river water temperature is strongly persistent and strongly seasonal, and
machine-learning studies routinely report large skill gains for it. Those gains
are almost always measured against naive persistence, climatology, or an
air--water regression, on random data splits, with each model scored on its own
set of predictable days. Such designs cannot separate learned river behaviour
from three cheaper explanations: seasonal damping, favourable key selection, and
spatial interpolation between neighbouring gauges. Here we score a constrained
deep predictor against damped persistence, gradient-boosted trees with and
without site identity, a global LSTM, and an information-matched plain causal
convolutional network on one common set of 249,072 station/\allowbreak date/\allowbreak horizon keys, at
120 U.S. gauges in 15 hydrologic regions over 2006--2020, with all preprocessing
fitted strictly backwards in time and whole regions, not random sites, held out.
Median station skill (one minus the RMSE ratio) against
persistence is +0.251 at a seven-day lead but +0.038 against damped persistence:
of the 0.560 °C reduction in station-median RMSE from persistence to the deep
model, 0.479 °C is delivered by damping alone. A gradient-boosted tree has the
lowest station-median RMSE at all three leads, and the information-matched plain
convolutional network is within 0.023 °C of the full architecture. Three-day
skill against persistence falls from +0.187 under a random held-site split to
+0.116 under whole-region holdout, at a mean nearest-training-gauge distance of
289 km. Reported skill for this variable is therefore largely a statement about
the reference model and the spatial partition rather than about the architecture.

% agujournal2025.cls provides a native Plain Language Summary environment, nested
% inside the abstract.  This replaces the \section*{Plain Language Summary}
% workaround the 2019 class required.  No abstract body may follow it.
\begin{plainlanguagesummary}
River temperature changes slowly from day to day, so
a forecast that repeats yesterday\textquotesingle s reading is already fairly accurate, and one
that also nudges it toward the usual value for the time of year is better still.
That makes it easy to publish a model that looks impressive without being useful.
We tested a river-temperature model at 120 U.S. gauges over fifteen years and
deliberately made the test hard: every model was scored on exactly the same days
and sites, none could see information from after the moment it was asked to
predict, and we held out entire river regions rather than scattered gauges.
Measured the usual way, our model looked strong. Measured against a simple
seasonal adjustment of yesterday\textquotesingle s reading, almost all of that advantage
disappeared, and a standard tree-based method was more accurate at every forecast
range. How well such a model is reported to perform depends mostly on what it is
compared against and on how the map is divided, not on how the model is built.
\end{plainlanguagesummary}
\end{abstract}

\noindent\textbf{Keywords:} river water temperature; benchmark design; strong baselines;
spatial holdout; clustered inference; conformal prediction.

\section{Introduction}\label{introduction}

Stream and river thermal regimes set dissolved-oxygen saturation, metabolic
rates, life-stage timing, and habitat suitability, and they respond to
atmospheric forcing, discharge, groundwater exchange, riparian shading, channel
geometry, and regulation
(\href{https://doi.org/10.1111/j.1365-2427.2006.01597.x}{Caissie, 2006}). Water
managers need daily-resolution temperature at gauged reaches to time releases,
anticipate thermal refuge loss, and interpret compliance monitoring, and the
supply of daily records from national monitoring networks has made the variable
one of the most heavily modelled targets in applied hydrological machine
learning.

The prevailing understanding in that literature is that deep sequence models have
substantially advanced daily river-temperature prediction. Nonlinear air--water
regressions (\href{https://doi.org/10.1029/98WR01877}{Mohseni et al., 1998}) and
hybrid air-temperature and discharge formulations
(\href{https://doi.org/10.1088/1748-9326/10/11/114011}{Toffolon and Piccolroaz, 2015};
\href{https://doi.org/10.1002/hyp.10913}{Piccolroaz et al., 2016}) were succeeded by
recurrent, multi-task, and physics-guided river-network architectures that report
large error reductions
(\href{https://doi.org/10.5194/hess-25-2951-2021}{Feigl et al., 2021};
\href{https://doi.org/10.1088/1748-9326/abd501}{Rahmani, Lawson, et al., 2021};
\href{https://doi.org/10.1137/1.9781611976700.69}{Jia et al., 2021};
\href{https://doi.org/10.1029/2021WR030138}{Sadler et al., 2022};
\href{https://doi.org/10.3389/frwa.2023.1184992}{Zwart et al., 2023}). The reported
gains are consistent enough that architectural sophistication is now the usual
explanation for accuracy in this problem.

Those reported gains are, however, extraordinarily heterogeneous, and the
heterogeneity does not organise itself by architecture. Reviews of the field note
that published accuracies are not comparable across studies because reference
models, key sets, and data-splitting conventions differ
(\href{https://doi.org/10.5194/hess-29-2521-2025}{Corona and Hogue, 2025}). The
dispersion is larger than the architectural differences it is supposed to
explain, and it survives within single panels: on the data used here, one fixed
model and one fixed set of prediction days produce a seven-day skill of +0.251 or
of +0.038 depending only on which reference model is placed in the denominator.
An architectural account cannot produce that spread, because the architecture
does not change between the two numbers.

The identification gap is a property of the designs, not of the literature\textquotesingle s
size. Three design choices are made almost universally, and each one absorbs an
alternative explanation into the reported number. First, skill is quoted against
naive persistence, climatology, or an air--water regression. Daily mean water
temperature has large thermal inertia, so relaxing the last observation toward a
seasonal climatology already removes most of the error a learned model can
remove; a reference that omits this hands the model a large part of its budget
before it learns anything. Second, models are scored on model-specific
complete-case sets, so a model can gain apparent accuracy by declining to predict
on hard days, and preprocessing statistics --- scaling constants, imputation fills,
climatological references, event thresholds, interval offsets --- are commonly
fitted on periods that include the evaluated interval, which transfers
information backwards in time without leaving a trace in a conventional split.
Third, spatial generalisation is reported from random held-site splits. When a
held-out gauge\textquotesingle s neighbours remain in training, the model reaches the held-out
regime through correlated forcing and spatially smooth representations; the
quantity measured is closer to interpolation than to transfer. None of these
choices is detectable from a reported RMSE, and each inflates it.

Here we hold the model and the data fixed and vary the three design choices
instead, so that their separate contributions to reported skill become
measurable. We assembled a frozen panel of 657,480 site-days from 120 stable U.S.
Geological Survey site numbers spanning 2006--2020, 34 states, and 15 two-digit
hydrologic unit (HUC2) regions, and evaluated one common set of
station/\allowbreak date/\allowbreak horizon keys at 1-, 3-, and 7-day leads. Four questions organise the
study. How much of a reported gain survives a strong reference? Which model class
is actually most accurate on identical keys? How much of a reported spatial
transfer survives whole-region holdout? And how far do those answers vary across
regions, seasons, and the width of a calibrated interval? The predictor whose
skill we decompose, ThermoRoute, is described in full in Section 3.1; its
architecture is the object under test rather than the contribution, and the
evaluation-period labels covering 2021--2023 are acquired exactly once, after the
model suite and analysis code are frozen and replayed.

\begin{center}\rule{0.5\linewidth}{0.5pt}\end{center}

\section{Data and design}\label{data-and-design}

\subsection{The frozen panel and the cohort it defines}\label{the-frozen-panel-and-the-cohort-it-defines}

The development record is a single derived panel,
\texttt{data\_\allowbreak{}usgs/\allowbreak{}panel\_\allowbreak{}usgs\_\allowbreak{}120v2.parquet}, bound by a manifest
(\texttt{data\_\allowbreak{}usgs/\allowbreak{}frozen\_\allowbreak{}panel\_\allowbreak{}v1.json}) to a station registry
(\texttt{data\_\allowbreak{}usgs/\allowbreak{}station\_\allowbreak{}registry\_\allowbreak{}v1.csv}). Each of 120 sites carries one row for
every calendar day from 2006-01-01 through 2020-12-31, giving 657,480 rows with
no gaps in the date index; missing observations are explicit nulls rather than
absent rows. Station identity is the zero-padded USGS \texttt{site\_\allowbreak{}no}.

The seven raw variables consumed by the models, in frozen order, are water
temperature (\texttt{WTEMP}), streamflow (\texttt{FLOW}), air temperature (\texttt{TEMP}),
precipitation (\texttt{PRCP}), a relative-humidity proxy (\texttt{RHMEAN}), Daymet
daylight-period mean incoming shortwave radiation in W m⁻² (\texttt{DH}), and wind speed
(\texttt{WDSP}). Two names are legacy and easily misread: \texttt{DH} is Daymet \texttt{srad}, the
incident shortwave flux averaged over the daylight period, not day length and not
a 24-hour mean; \texttt{RHMEAN} is a reproducible vapour-pressure proxy evaluated at the
tmax/tmin midpoint, not a direct daily-mean humidity observation. Gage height
(\texttt{WLEVEL}) is retained as raw evidence but is not a model input. Water
temperature and discharge come from USGS NWIS daily values
(\href{https://doi.org/10.5066/F7P55KJN}{U.S. Geological Survey, 2024}),
meteorological fields from Daymet V4
(\href{https://doi.org/10.3334/ORNLDAAC/2129}{Thornton et al., 2022}), and wind from
gridMET (\href{https://doi.org/10.1002/joc.3413}{Abatzoglou, 2013}).

The cohort was not sampled at random from U.S. rivers, and nothing here treats it
as though it were. Candidate selection began from a discovery query over stream
sites in the states represented by the panel and rejected 1,345 of 1,465
candidate stations: 950 lacked joint water-temperature and discharge
availability, 376 failed a full-period coverage threshold, and 19 failed a
2019--2020 coverage threshold. The retained cohort is therefore
availability-enriched, and the 2019--2020 interval participated in its
construction, so that interval is development data in the strict sense and is
exploratory throughout this manuscript.

Three properties of the retained cohort bear directly on the analysis. Observed
water temperature is absent on about 15.8\% of panel rows, with a maximum
contiguous gap of 2,404 days, and discharge on about 2.8\%, with a maximum gap of
1,369 days; the meteorological fields are missing on about 0.07\% of rows, in a
structured rather than random way, because all 120 sites lack the provider
calendar\textquotesingle s final day in the leap years 2008, 2012, 2016, and 2020. Among retained
stations, 38 share a repeated hydrologic unit code with another retained station
and 19 have another retained station within 10 km, so station-level independence
is not tenable. Two stations contain 2,059 rows of signed negative discharge
(minimum −121 cfs), retained with their sign rather than clipped.

\subsection{Temporal roles and the common key set}\label{temporal-roles-and-the-common-key-set}

Temporal roles are fixed before any model is fitted, and a training sample is
admitted only if its issue date and every one of its target dates fall inside the
same partition.

{\small
\setlength{\tabcolsep}{3pt}
\renewcommand{\arraystretch}{1.12}
\noindent
\begin{tabularx}{\dimexpr\linewidth-24pt\relax}{@{}lrrr>{\raggedright\arraybackslash}X@{}}
\toprule\noalign{}
Role & Dates & Rows & Observed \texttt{WTEMP} & Permitted use \\
\midrule\noalign{}
Training & 2006--2015 & 438,240 & 341,646 & fit preprocessing and models \\
Validation & 2016--2017 & 87,720 & 84,074 & choose frozen settings \\
Calibration & 2018 & 43,800 & 42,279 & conformal and Platt fits; final year of the 2006--2018 seasonal event reference \\
Development evaluation & 2019--2020 & 87,720 & 85,621 & previously inspected; exploratory diagnosis only \\
One-time evaluation & 2021--2023 & --- & sealed & opened once, after all freezes (Section 3.7) \\
\bottomrule\noalign{}
\end{tabularx}
}

On the development-evaluation partition, the intersection of admissible keys
across all primary models comprises 249,072 station/\allowbreak date/\allowbreak horizon keys, evenly
distributed as 83,024 keys per lead. Every development comparison in Section 4
uses exactly this key set. No model-specific complete-case set is permitted, so a
model cannot gain apparent accuracy by declining to predict on hard keys.

\subsection{Evaluation-period inputs and the information boundary}\label{evaluation-period-inputs-and-the-information-boundary}

The evaluation interval is 2021-01-01 through 2023-12-31 for the same 120-site
cohort, and historical Daymet and gridMET covariates for that interval are
retrieved and archived before any outcome is accessed. Meteorology is represented
at each station coordinate and is not aggregated over upstream catchments, which
is a real limitation for large basins (Section 6.2).

The primary information set uses provider values dated no later than each
historical issue date and consumes no horizon-specific future weather forecast.
This is a date-indexed retrospective hindcast, and not a re-execution of what a
forecaster could have run on the day (Section 6.1): the
as-issued provisional vintage of a gridded product cannot be reconstructed after
the fact, so archiving requests, responses, timestamps, and checksums freezes the
dataset actually evaluated without proving that identical values were available
operationally at the time. NWIS dates denote site-local finalized daily values
while Daymet and gridMET use provider-specific calendar-day definitions; no
subdaily day-boundary harmonization is claimed. A separate outcome-free bridge,
described in the Supporting Information, re-fetched 2018--2020 covariates with the
parser used for the evaluation-period acquisition and records
\texttt{PASS\_\allowbreak{}EXACT\_\allowbreak{}PRODUCT\_\allowbreak{}BRIDGE} on the exact site/date registry.

A restricted external cohort is defined after the model suite is committed: a
fixed 34-state metadata query identifies stream sites advertising daily
water-temperature capability, without requesting values, coverage dates, or event
rates, and a deterministic seed selects 30 stations whose site identifiers do not
occur in the development registry. This is site-identifier disjointness only, and
those models still consume each target site\textquotesingle s observed history (Section 3.4).

\begin{center}\rule{0.5\linewidth}{0.5pt}\end{center}

\section{Methods}\label{methods}

\subsection{The predictor}\label{the-predictor}

ThermoRoute predicts at issue time \emph{t}, station \emph{i}, and lead \emph{h} by adding a
bounded learned correction to a damped-persistence anchor. One model is fitted
for all three leads jointly: the horizon set (1, 3, 7 days) is an input dimension
and every head emits a lead-indexed vector in a single forward pass, so no lead
has its own weights.

\textbf{Anchor.} A frozen seasonal climatology \texttt{c} and a per-station decay \texttt{φ\_\allowbreak{}i} give

\begin{quote}
\textbf{(1)} \texttt{A\_\allowbreak{}\{i,t+h\}\ \allowbreak{}=\ \allowbreak{}c\_\allowbreak{}\{i,t+h\}\ \allowbreak{}+\ \allowbreak{}φ\_\allowbreak{}i\^{}h\ \allowbreak{}·\ \allowbreak{}(y\_\allowbreak{}\{i,t\}\ \allowbreak{}−\ \allowbreak{}c\_\allowbreak{}\{i,t\})}
\end{quote}

where \texttt{y\_\allowbreak{}\{i,t\}} is the last observed daily mean water temperature. Each \texttt{φ\_\allowbreak{}i} is
estimated by no-intercept least squares on consecutive-day anomaly pairs from the
training period only, requires at least 30 pairs, and is clipped to {[}0, 0.999{]}.
This object, fitted identically, is also the damped-persistence reference.

\textbf{Learned relaxation proposal.} A station-, flow-, and season-conditioned
relaxation rate replaces the fixed decay in a separate proposal path:

\begin{quote}
\textbf{(2)} \texttt{κ\_\allowbreak{}\{i,t\}\ \allowbreak{}=\ \allowbreak{}σ(b\ \allowbreak{}+\ \allowbreak{}b\_\allowbreak{}i\ \allowbreak{}+\ \allowbreak{}β\_\allowbreak{}q\ \allowbreak{}q\_\allowbreak{}\{i,t\}\ \allowbreak{}+\ \allowbreak{}β\_\allowbreak{}s\ \allowbreak{}s\_\allowbreak{}t)}, clipped to {[}10⁻³, 0.999{]}

\textbf{(3)} \texttt{P\_\allowbreak{}\{i,t+h\}\ \allowbreak{}=\ \allowbreak{}c\_\allowbreak{}\{i,t+h\}\ \allowbreak{}+\ \allowbreak{}e\_\allowbreak{}\{i,t\}\ \allowbreak{}+\ \allowbreak{}(1\ \allowbreak{}−\ \allowbreak{}κ\_\allowbreak{}\{i,t\})\^{}h\ \allowbreak{}·\ \allowbreak{}(a\_\allowbreak{}\{i,t\}\ \allowbreak{}−\ \allowbreak{}e\_\allowbreak{}\{i,t\})}
\end{quote}

with \texttt{a\_\allowbreak{}\{i,t\}\ \allowbreak{}=\ \allowbreak{}y\_\allowbreak{}\{i,t\}\ \allowbreak{}−\ \allowbreak{}c\_\allowbreak{}\{i,t\}} the current anomaly, \texttt{e\_\allowbreak{}\{i,t\}} a learned
equilibrium anomaly formed from the standardized forcings and a station term,
\texttt{q} the standardized log-discharge, and \texttt{s\_\allowbreak{}t} a two-component season encoding.
The bias \texttt{b} is initialised at −2.94, so the proposal begins near damped
persistence. The fitted relaxation rate is a statistical quantity only: it is not
a heat-transfer coefficient, and no physical interpretation of its value is
offered anywhere in this paper.

\textbf{Router.} A sparse, horizon-conditioned router scores all 7 variables at lags 0
to 14 --- 105 (variable, lag) pairs --- with a scaled dot product between a
horizon-specific query and a per-pair key, normalised by sparsemax
(\href{https://proceedings.mlr.press/v48/martins16.html}{Martins and Astudillo, 2016}):

\begin{quote}
\textbf{(4)} \texttt{w\_\allowbreak{}h\ \allowbreak{}=\ \allowbreak{}sparsemax(\ \allowbreak{}⟨q\_\allowbreak{}h\ \allowbreak{}+\ \allowbreak{}ctx,\ \allowbreak{}k\_\allowbreak{}\{v,ℓ\}⟩\ \allowbreak{}/\allowbreak{}\ \allowbreak{}√d\ \allowbreak{})}, over \texttt{v\ \allowbreak{}∈\ \allowbreak{}1…7}, \texttt{ℓ\ \allowbreak{}∈\ \allowbreak{}0…14}
\end{quote}

The router is an allocation mechanism inside the predictor, not a map of
connected reaches; it receives no verified graph or topology input.

\textbf{Temporal encoder.} A strictly left-looking temporal convolutional network
(\href{https://arxiv.org/abs/1803.01271}{Bai et al., 2018}) of \texttt{B\ \allowbreak{}=\ \allowbreak{}2} residual blocks
with kernel \texttt{k\ \allowbreak{}=\ \allowbreak{}3}, dilations 1 and 2, and width 40 encodes recent history. Left
padding of \texttt{(k\ \allowbreak{}−\ \allowbreak{}1)·2\^{}\{b\}} makes the output at position \emph{t} a function of
positions ≤ \emph{t} by construction. The receptive field is

\begin{quote}
\textbf{(5)} \texttt{R\ \allowbreak{}=\ \allowbreak{}1\ \allowbreak{}+\ \allowbreak{}(k\ \allowbreak{}−\ \allowbreak{}1)(2\^{}B\ \allowbreak{}−\ \allowbreak{}1)\ \allowbreak{}=\ \allowbreak{}7} steps
\end{quote}

so, together with the router\textquotesingle s oldest usable value at lag 14, no input older than
lag 14 can affect the output. The sequence builder supplies a 32-day tensor, but
that is a construction buffer, not an effective-memory claim. Throughout this
paper "causal" describes time ordering only and never causal inference.

\textbf{Regime mixture.} Three experts are combined by a soft gate
(\href{https://arxiv.org/abs/1701.06538}{Shazeer et al., 2017}):

\begin{quote}
\textbf{(6)} \texttt{r\ \allowbreak{}=\ \allowbreak{}Σ\_\allowbreak{}\{k=1\}\^{}\{3\}\ \allowbreak{}π\_\allowbreak{}k\ \allowbreak{}·\ \allowbreak{}E\_\allowbreak{}k(·)}, \texttt{π\ \allowbreak{}=\ \allowbreak{}softmax(g(·))}
\end{quote}

\textbf{Bounded residual.} The learned point correction is squashed and rescaled, so
its deviation from the anchor is algebraically bounded:

\begin{quote}
\textbf{(7)} \texttt{ŷ\_\allowbreak{}\{i,t+h\}\ \allowbreak{}=\ \allowbreak{}A\_\allowbreak{}\{i,t+h\}\ \allowbreak{}+\ \allowbreak{}δ\ \allowbreak{}·\ \allowbreak{}tanh(\ \allowbreak{}z\_\allowbreak{}\{i,t+h\}\ \allowbreak{}/\allowbreak{}\ \allowbreak{}δ\ \allowbreak{})}, \texttt{δ\ \allowbreak{}=\ \allowbreak{}1.0\ \allowbreak{}°C}
\end{quote}

giving \texttt{\textbar{}ŷ\ \allowbreak{}−\ \allowbreak{}A\textbar{}\ \allowbreak{}\textless{}\ \allowbreak{}δ} identically. The bound is relative to the named anchor: it
limits the point output\textquotesingle s deviation from damped persistence and bounds neither
absolute error, nor event-tail error, nor interval width, nor behaviour after a
distribution shift. Because the 2019--2020 partition was inspected during
development, \texttt{δ\ \allowbreak{}=\ \allowbreak{}1.0\ \allowbreak{}°C} is a development-selected algebraic point bound and is
not a selection from an untouched holdout.

\textbf{Loss.} With \texttt{y} the observed target, \texttt{A} the anchor, quantile levels
τ ∈ \{0.05, 0.50, 0.95\}, \texttt{Q\_\allowbreak{}τ} the corresponding quantile head, and \texttt{ρ\_\allowbreak{}τ} the
pinball loss,

\begin{quote}
\textbf{(8)} \texttt{L\ \allowbreak{}=\ \allowbreak{}MSE(y,\ \allowbreak{}ŷ)\ \allowbreak{}+\ \allowbreak{}Σ\_\allowbreak{}τ\ \allowbreak{}ρ\_\allowbreak{}τ(y,\ \allowbreak{}Q\_\allowbreak{}τ)\ \allowbreak{}+\ \allowbreak{}λ\_\allowbreak{}e\ \allowbreak{}·\ \allowbreak{}BCE(exceedance)\ \allowbreak{}+\ \allowbreak{}λ\_\allowbreak{}c\ \allowbreak{}·\ \allowbreak{}C\ \allowbreak{}+\ \allowbreak{}λ\_\allowbreak{}r\ \allowbreak{}·\ \allowbreak{}‖ŷ\ \allowbreak{}−\ \allowbreak{}A‖₁}
\end{quote}

with \texttt{λ\_\allowbreak{}e\ \allowbreak{}=\ \allowbreak{}0.3}, \texttt{λ\_\allowbreak{}r\ \allowbreak{}=\ \allowbreak{}10⁻²}, \texttt{λ\_\allowbreak{}c\ \allowbreak{}=\ \allowbreak{}1.0}, and a temperature loss scale of 1.0.
The point and pinball terms carry unit weight. The term \texttt{C} penalises quantile
non-monotonicity and is identically zero under the head parameterisation used
here, so it contributes no gradient; it is retained because it was declared
before fitting.

\textbf{Fitting.} All neural members are fitted with AdamW
(Kingma and Ba, 2015; Loshchilov and Hutter, 2019) at learning rate 2 × 10⁻³,
weight decay 10⁻⁴, batch size 1,536, gradient-norm clipping at 1.0, a maximum of
80 epochs with early stopping at patience 12 on station-macro validation RMSE,
a plateau scheduler halving the rate at patience 4, and equal-station fixed-size
bootstrap sampling; five seeds (0--4) are fitted and averaged with equal weights.
The canonical configuration has \textbf{38,505 trainable parameters}. All members were
fitted on CPU only (Intel Xeon Gold 6430, one PyTorch intra-op and one inter-op
thread), and every seed stopped early, at epochs 18, 15, 14, 18, and 26. Measured
end-to-end training wall-clock time is \texttt{{[}TRAINING\ \allowbreak{}WALL-CLOCK\ \allowbreak{}TO\ \allowbreak{}BE\ \allowbreak{}RECORDED{]}} and
inference throughput is \texttt{{[}INFERENCE\ \allowbreak{}COST\ \allowbreak{}TO\ \allowbreak{}BE\ \allowbreak{}RECORDED{]}}; neither is present in
the run records, and neither is estimated here.

\subsection{The reference set}\label{the-reference-set}

The composition of the reference set is the most consequential methodological
choice in this study, because it is what the reported gain is a gain over. Each
member removes one alternative explanation for apparent skill.

\textbf{Persistence} (\texttt{ŷ\_\allowbreak{}\{t+h\}\ \allowbreak{}=\ \allowbreak{}y\_\allowbreak{}t}) quantifies the raw memory of the series.
\textbf{Damped persistence}, equation (1) fitted on the training period, absorbs the
two effects a learned model most easily rediscovers --- inertia and seasonality ---
and is the reference against which the primary comparisons are stated.
\textbf{Seasonal climatology} isolates the seasonal cycle alone. \textbf{LightGBM}
(\href{https://proceedings.neurips.cc/paper/2017/hash/6449f44a102fde848669bdd9eb6b76fa-Abstract.html}{Ke et al., 2017})
is the principal learned reference and is run both as a global model receiving
stable site identity as a categorical feature and as a per-station variant that
isolates the value of pooling. A \textbf{global LSTM} with a station embedding
represents the deep sequence family that has produced the strongest recent
results for this variable
(\href{https://doi.org/10.1088/1748-9326/abd501}{Rahmani, Lawson, et al., 2021};
\href{https://doi.org/10.3389/frwa.2023.1184992}{Zwart et al., 2023}). A \textbf{plain
causal temporal convolutional network} is the information-matched deep
reference: it receives the same outcome-free issue-time history, station
identity, climatology, frozen damped anchor, standardized environmental context,
and calendar and regime-gate tensors as the full model, is matched to within 2\%
of its parameter count (38,346 against 38,505) and to the same optimiser-step
budget, and predicts an unrestricted residual around the same anchor while
omitting the relaxation proposal, router, mixture, and residual bound. Its
companion, a plain multilayer perceptron at 38,860 parameters, isolates the value
of sequence structure. Neither ever reads the target tensor, the target-date
field, or the disabled water-level channel.

The learned references are given a genuine chance to win. The global LSTM uses
the same development splits, loss, station-balanced sampling, epoch and patience
limits, history length, site identity, and five-seed budget as ThermoRoute, its
predeclared validation-only grid may expose the same frozen climatology,
damped-anchor, and season features, and it is eligible for the same calibration;
it lacks only the bounded-residual constraint and the lag router. LightGBM
selects among four predeclared settings separately by lead on the 2016--2017
partition using station-macro RMSE alone; the LSTM selects among three
predeclared architectures with seed 0 before fitting five members. Tuning budgets
are documented but \emph{not} equalized across model classes, which is a real
asymmetry: any statement about LightGBM here is scoped to this frozen
four-candidate procedure and this feature schema, not to gradient boosting in
general. Seven one-factor architecture controls (\texttt{DampedPriorOnly},
\texttt{TR-noDynamicPrior}, \texttt{TR-fixedKappa}, \texttt{TR-noRouter}, \texttt{TR-noMoE}, \texttt{TR-noTCN},
\texttt{TR-unbounded}) are fitted with the same five seeds and paired within seed on
identical keys before ensemble averaging; they are deletion and intervention
sensitivities and do not prove component necessity or identify a mechanism.

An air2stream-style hybrid reference was declared but is recorded as \texttt{NOT\_\allowbreak{}RUN},
and the available implementation is an unofficial style reference rather than the
published code or a validated reproduction of it. It is therefore excluded from
every claim in this paper; its status, provenance, and what a defensible
comparison would require are reported in the Supporting Information rather than
carried as an unfilled column.

\subsection{Leakage control}\label{leakage-control}

Leakage control here is a set of mechanisms that can be checked against the code
and the artifacts, rather than a statement of intent.

\emph{Issue-time separation of covariates.} Every predictor consumed by a primary
model carries a date no later than the issue date, and no horizon-specific future
weather field enters any primary model.

\emph{Non-anticipating sequence encoding.} Equation (5) makes the information horizon
bounded and auditable by construction rather than by convention.

\emph{Strictly backward-fitted statistics.} Standardization constants, imputation
fills, the seasonal climatology, the damped-persistence rate, the station-level
q90 event thresholds (2006--2015), the seasonal event reference (2006--2018), the
conformal offsets (2018), and the Platt calibrators (2018) are each fitted on
data strictly preceding the interval on which they are applied. This closes the
most common backward-information channel in hydrological benchmarking, which is
not the model but the preprocessing.

\emph{Explicit admissibility.} A forecast key is admissible only when issue-date water
temperature is genuinely observed and a 32-day history can be constructed; other
history cells are filled with training-only seasonal medians and retain explicit
missingness masks. We record one honest gap: the issue-date auxiliary quantities
used by the relaxation proposal and the regime gate are computed from the same
training-only imputed panel but do not each carry a separate auxiliary-path mask,
so the output is not guaranteed invariant to the chosen fill values.

\emph{Replay under isolation.} Before any label acquisition is authorized, a fresh
isolated interpreter, denied network access, child processes, repository writes,
and reads from the evaluation namespaces, reloads every trained member and
reproduces the validation, calibration, and 2019--2020 keys and values from the
frozen inputs, independently recomputing the thresholds and refitting the
calibration parameters; a prediction file that is merely self-consistent is
rejected. The isolation model, the byte bindings, and the publication and crash
semantics of that replay are specified in the Supporting Information.

\subsection{Two spatial partitions}\label{two-spatial-partitions}

Random held-site splits are the usual way to report spatial generalisation in
this literature, and they leak. If a held-out gauge\textquotesingle s neighbours remain in
training, the model reaches the held-out site\textquotesingle s thermal regime through correlated
forcing, shared preprocessing statistics, and spatially smooth learned
representations. We therefore report two arms with different leakage properties
and read the difference between them as the measurement of interest.

The \textbf{random held-site warm-start} arm uses four folds in which held-site
histories still contribute to global panel preprocessing. This is the arm most
comparable to published random-split results, and precisely for that reason it is
not the arm we treat as the spatial test.

The \textbf{held-region gauged transfer} arm packs the 15 HUC2 groups into four folds
of {[}30, 30, 31, 29{]} stations and holds out \emph{whole regions}, so no gauge from a
held-out region appears in training. Station-agnostic ThermoRoute and a global
LightGBM are each trained on the in-fold regions and forecast the held-out
region\textquotesingle s stations, with all climatology, scaling, and damped-rate parameters
pooled from in-fold stations only. The mean distance from a held-out station to
its nearest training gauge is 289 km.

The label on this arm must be exact. Held-out sites still supply their own
observed water temperature through the issue date, both as the persistence anchor
and as the sequence history. This is \textbf{gauged} transfer to a region whose gauges
were not used in fitting; it is \textbf{not} prediction at a site with no
observational record, which is a different problem and a different evaluation
(\href{https://doi.org/10.3390/w14071032}{Weierbach et al., 2022};
\href{https://doi.org/10.1002/hyp.14400}{Rahmani, Shen, et al., 2021}). The same
caveat applies to the 30-site external cohort of Section 2.3.

\subsection{Conformal intervals}\label{conformal-intervals}

Predictive uncertainty in hydrological modelling has most often been represented
through Bayesian or likelihood-based treatments of parameter and input error
(\href{https://doi.org/10.1002/hyp.3360060305}{Beven and Binley, 1992};
\href{https://doi.org/10.1029/2005WR004368}{Kavetski et al., 2006a},
\href{https://doi.org/10.1029/2005WR004376}{2006b}). We use a distribution-free
calibration step instead, because the quantity required here is an empirical
coverage statement about a fixed set of forecast keys rather than a posterior
over model parameters. Intervals come from split conformalized quantile
regression
(\href{https://papers.nips.cc/paper/2019/hash/5103c3584b063c431bd1268e9b5e76fb-Abstract.html}{Romano et al., 2019};
\href{https://doi.org/10.1007/b106715}{Vovk et al., 2005}): the learned models emit a
separate mean-squared-error point head and three pinball-trained quantile heads,
members are averaged with equal weights, and the offset is fitted on the 2018
calibration year only. The construction is deliberately one-sided --- the deployed
offset is \texttt{qhat\_\allowbreak{}plus\ \allowbreak{}=\ \allowbreak{}max(raw\_\allowbreak{}qhat,\ \allowbreak{}0)} and the delivered interval is
\texttt{{[}q05\ \allowbreak{}−\ \allowbreak{}qhat\_\allowbreak{}plus,\ \allowbreak{}q95\ \allowbreak{}+\ \allowbreak{}qhat\_\allowbreak{}plus{]}} with the median unchanged --- so calibration
can widen a nominal interval but never shrink it, and a non-finite, crossed, or
empty final interval fails closed before any evaluation label is read.

What this guarantees is worth stating precisely, because conformal prediction is
frequently over-claimed in applied work. Split conformal provides finite-sample
\emph{marginal} coverage under exchangeability of calibration and evaluation
nonconformity scores. Here the calibration scores come from 2018 and the
evaluation scores from a later, disjoint period at the same sites, so
exchangeability fails by construction: there is temporal drift, and the scores
are dependent within station and within region. We therefore report achieved
coverage as an \emph{empirical marginal diagnostic}, make no finite-sample guarantee,
and make no conditional-coverage claim of any kind. Two pre-specified
sensitivities probe that boundary: a block-maximum variant that calibrates on the
maximum score within groups of consecutive retained calibration rows, following
the standard practice of resampling blocks rather than rows when the series is
dependent (\href{https://doi.org/10.1214/aos/1176347265}{Künsch, 1989}), and an
idealized delayed adaptive-conformal variant using each forecast\textquotesingle s target date as
a feedback-arrival proxy. Neither replays real feedback availability, because the
inputs contain no verified observation-publication timestamp, revision history,
or reporting latency. The equal-weight three-quantile pinball summary is not
CRPS; the two are distinct members of the family of proper scoring rules
(\href{https://doi.org/10.1198/016214506000001437}{Gneiting and Raftery, 2007}). An
event head separately reports exceedance of each station\textquotesingle s 2006--2015 q90 water
temperature, calibrated by one Platt map per lead fitted on 2018 only; that
threshold is a statistical tail diagnostic local to each station, with no
biological, ecological, or regulatory interpretation.

\subsection{Estimand, metrics, comparison set, and why this study is descriptive}\label{estimand-metrics-comparison-set-and-why-this-study-is-descriptive}

\textbf{Two differently signed quantities are reported, and they are never combined.}
The paired effect is

\begin{quote}
\textbf{(9)} \texttt{ΔRMSE\ \allowbreak{}=\ \allowbreak{}RMSE(candidate)\ \allowbreak{}−\ \allowbreak{}RMSE(reference)}, in °C, \textbf{negative} favours the candidate
\end{quote}

and the skill score is

\begin{quote}
\textbf{(10)} \texttt{skill\ \allowbreak{}=\ \allowbreak{}1\ \allowbreak{}−\ \allowbreak{}RMSE(candidate)/\allowbreak{}RMSE(reference)}, dimensionless, \textbf{positive} favours the candidate
\end{quote}

Equation (9) is the formal estimand and carries a °C unit everywhere it appears;
equation (10) is a ratio and appears as a bare signed number. Every value in
Section 4 is labelled with which of the two it is. A positive ΔRMSE and a
positive skill score mean opposite things, and no figure axis, table column, or
sentence in this paper mixes them.

The sampling unit is the station. For each lead, unweighted RMSE is computed on
the common daily keys separately for each reportable station, and the primary
effect is the unweighted median across stations of equation (9) with ThermoRoute
as candidate; a station/lead cell is reportable only with at least 100 valid
paired targets, so daily rows do not determine between-station weight. We report
RMSE and paired RMSE differences rather than an efficiency-type criterion such as
the Nash--Sutcliffe efficiency or its decomposition-based successors
(\href{https://doi.org/10.1016/0022-1694(70)90255-6}{Nash and Sutcliffe, 1970};
\href{https://doi.org/10.1016/j.jhydrol.2009.08.003}{Gupta et al., 2009}), because
those criteria normalise by a station-specific variance, which would make a
paired between-model difference on identical keys harder to read.

The frozen comparison set contains exactly five rows:

{\small
\setlength{\tabcolsep}{3pt}
\renewcommand{\arraystretch}{1.12}
\noindent
\begin{tabularx}{\dimexpr\linewidth-24pt\relax}{@{}l>{\raggedright\arraybackslash}Xr>{\raggedright\arraybackslash}X@{}}
\toprule\noalign{}
\# & Comparison & Lead & Predeclared margin \\
\midrule\noalign{}
1 & ThermoRoute vs. damped persistence & 1 d & 0.00 °C \\
2 & ThermoRoute vs. damped persistence & 3 d & 0.00 °C \\
3 & ThermoRoute vs. damped persistence & 7 d & 0.00 °C \\
4 & ThermoRoute vs. LightGBM & 3 d & +0.05 °C numerical ceiling \\
5 & ThermoRoute vs. LightGBM & 7 d & +0.05 °C numerical ceiling \\
\bottomrule\noalign{}
\end{tabularx}
}

The +0.05 °C ceiling was fixed before outcomes as a strict numerical limit on
allowable degradation. It is not derived from sensor precision, biological
response, a water-quality standard, or an elicited stakeholder utility, and it
carries no ecological or regulatory importance.

Two uncertainty procedures accompany each row, both clustered at HUC2: a
one-sided p-value from exact enumeration of all 2\^{}K whole-cluster sign vectors,
applying one common sign to every station effect within a cluster, and a
10,000-draw whole-HUC2 cluster bootstrap percentile interval for the median
station effect, with Holm adjustment
(\href{https://www.jstor.org/stable/4615733}{Holm, 1979}) over exactly these five
p-values. We do not use the standard tests of equal predictive accuracy for
forecast series
(\href{https://doi.org/10.1080/07350015.1995.10524599}{Diebold and Mariano, 1995};
\href{https://doi.org/10.1016/S0169-2070(96)00719-4}{Harvey et al., 1997}), because
their asymptotics are stated for a single loss-differential series and do not
address dependence across the 120 gauges, which is the dominant dependence here.

\textbf{Those procedures cannot carry inferential weight in this study, for three
independent and permanent reasons.} Eligibility is a fixed property of the
frozen gate rather than a quantity recomputed from the cohort, and the opening
refuses to run unless the gate has failed closed. The reasons follow in order of
how binding they are.

\emph{First, the two structural assumptions on which both procedures rest are recorded
as unmet.} Whole-cluster resampling assumes independent, exchangeable sampling of
HUC2 regions, and whole-cluster sign flipping assumes joint sign symmetry of each
complete cluster effect vector around the tested margin. The frozen gate records
both as permanently unmet, for reasons that are design facts rather than
measurements: the regions were not probability-sampled, and no randomized sign
assignment exists anywhere in the design. This is not a threshold a larger or
differently partitioned cohort would clear. It is a statement about how the
cohort came into being --- an availability filter over gauges in 34 states --- and it
would hold with 30 regions as surely as with 15.

\emph{Second, the null-simulation component was never implemented.} The gate requires
a simulated-null calibration of the clustered procedures before they may be read
inferentially. That component does not exist in the analysis code; it returns a
fixed failure on every run and appends its own blocking reason. It is a declared
requirement that was not built, disclosed rather than quietly dropped.

\emph{Third, and least binding, the cohort has too few clusters.} The gate requires at
least 30 reportable clusters, an effective-cluster fraction of at least 0.75, and
a largest-cluster share below 0.25. Before any reportability attrition the 15
HUC2 groups contain between 2 and 26 stations, the largest holds 21.7\% of
stations, and the inverse-Herfindahl effective cluster count is 9.54, an effective
fraction of about 0.636; two of the three components fail. We state plainly what
this means: \textbf{the threshold is unreachable under a HUC2 partition for any U.S.
cohort whatsoever}, because the Watershed Boundary Dataset defines about 21 HUC2
regions in total and the gate asks for at least 30. That is a specification
error, made when the amendment was written and identified before any evaluation
label was accessed --- which is the only reason it can be reported here rather than
discovered by a reader.

Disclosing the error is a better response than defending the threshold, and a far
better response than subdividing the map until it passes. Subdividing \emph{would}
pass: recomputed with the gate\textquotesingle s own code on the frozen registry, HUC2 gives 15
clusters at an effective fraction of 0.636, HUC4 gives 64 at 0.507, HUC6 gives 75
at 0.485, and HUC8 gives 95 at 0.758, clearing all three thresholds. We did not
adopt HUC8. Adjacent HUC8 units on the same river are plainly not independent, so
passing that way would satisfy the arithmetic while defeating its purpose; and it
would change nothing, because eligibility is fixed by the first two reasons, so a
HUC8 partition would clear the cluster gate and return the same verdict. Choosing
the cluster definition after seeing which definition helps is precisely the
failure mode the gate exists to prevent.

Because the frozen Route-A cohort contains at most 15 HUC2 groups (inverse-Herfindahl effective cluster count ≈ 9.54) and therefore fails the outcome-free inference gate requiring n\_clusters ≥ 30, effective\_cluster\_fraction ≥ 0.75, and largest\_cluster\_share \textless{} 0.25, the formal five confirmatory comparisons are permanently restricted to fixed-cohort descriptive effects under DESCRIPTIVE\_ONLY\_INFERENCE\_GATE\_FAILED for confirmatory claim eligibility; exploratory and secondary analyses remain reportable as such and are not erased by the gate. Whole-HUC2 bootstrap intervals, exact sign-flip p-values, and Holm adjustments are assumption-conditional sensitivities only and must not be written as confirmatory decision evidence; after opening they must not be rewritten as superiority, non-inferiority, equivalence, parity, or national / U.S.-river generalization.

Every one of the five rows is still rendered exactly once, including adverse,
conflicting, or non-reportable outcomes, with the effect estimate, the win rate,
the station and cluster counts, and the interval all shown. Secondary and
exploratory analyses --- a qualifier-restricted target sensitivity, equal-region
and leave-one-region influence summaries, station-balanced probability
diagnostics, architecture controls, and synthetic perturbations --- are labelled as
such and cannot promote or replace a formal row. Relative economic value in the
cost--loss sense (\href{https://doi.org/10.1002/qj.49712656313}{Richardson, 2000}) is
recorded as \texttt{REV\_\allowbreak{}NOT\_\allowbreak{}EVALUATED\_\allowbreak{}NO\_\allowbreak{}PREDECLARED\_\allowbreak{}COST\_\allowbreak{}LOSS\_\allowbreak{}RATIOS}, because no
cost--loss ratios, management costs, actions, or stakeholder utility model were
frozen. A predeclared coverage and calendar-balance audit accompanies each formal
row with eight fixed descriptive candidates --- equal weighting of the 12
year-by-season cells, leave-one-year for each evaluation year, and
leave-one-season for DJF, MAM, JJA, and SON --- of which the most adverse
candidate-minus-reference effect is taken as the deterministic unfavorable value,
with ties broken by frozen candidate order. That audit is restricted to keys
where both issue-date and target-date water temperature are retained finite
observations; it does not impute unavailable outcomes, test a missing-at-random
assumption, or estimate all-calendar-day performance, and a favourable
sensitivity can never rescue an unfavourable primary row.

\subsection{The one-time 2021--2023 evaluation}\label{the-one-time-20212023-evaluation}

The evaluation labels are acquired exactly once, and everything capable of being
tuned is frozen and independently verified before that acquisition is permitted.
Those labels may not be read by model-selection, feature-selection,
threshold-selection, calibration, or station-inclusion code, and enforcement is
ordered so that the order is itself a gate. The protocol and its seal, the claim
renderer and validator, the source tree, the dependency locks, and all model
members are frozen and exactly replayed under the isolation of Section 3.3; the
model-suite registry is then committed \emph{while} the external candidate registry
and the evaluation-period predictor artifacts are still absent from the
repository; and only then may metadata-only candidate discovery and retrospective
predictor acquisition run. Authorization replays the resulting version-control
ancestry and content hashes to verify that the model-suite commit genuinely
precedes them. Any source change after the model freeze invalidates
authorization, and failure of the chronology gate does not produce a weaker
result: it demotes the exercise to retrospective exploration.

The acquisition itself is narrow and irreversible. A raw-only child process using
the USGS \texttt{dataretrieval} client
(\href{https://doi.org/10.5066/P94I5TX3}{Hodson and Hariharan, 2023}) is restricted to
daily mean water temperature (\texttt{00010}, °C), discharge (\texttt{00060}, cfs), and
raw-only gage height (\texttt{00065}, ft), and records exact request and response bytes,
qualifiers, timestamps, and content hashes. Where two or more finite series exist
for a parameter and date, the cell is marked missing with an explicit conflict
code rather than averaged or selected, and \textbf{no site may be replaced after
outcomes are accessed}. The primary analysis uses every finite parsed value
regardless of approval qualifier, so qualifiers can never select or remove a
station, model, date, or forecast key; a separately labelled sensitivity, fixed
before opening, retains target keys only when the target qualifier has the exact
token set \texttt{\{A\}}. The acquisition is transactional, with honest-owner crash and
replay guards rather than protection against a malicious owner or a
same-privilege adversary; those details are in the Supporting Information.

The analysis protocol was written and frozen before any post-2020 outcome was
requested or inspected, and fixes the cohort hashes, the temporal roles, the
model registry, the estimand, the five-row comparison set and its margins, the
clustered procedures, the multiplicity correction, the probabilistic and event
contracts, the outcome quality-control policy, and the requirement that every row
be reported exactly once regardless of direction; two later amendments and a
probability-metric erratum are themselves outcome-free and separately sealed. We
state the limit of that evidence plainly: the seals are repository-internal and
assume an honest owner, with no external timestamp, public registration service,
or independent custodian.

\begin{center}\rule{0.5\linewidth}{0.5pt}\end{center}

\section{Results}\label{results}

Sections 4.1--4.5 report development-period quantities. \textbf{All of them are
exploratory.} The 2019--2020 partition participated in cohort construction and
has informed model and narrative development; it is not an independent test, and
no number in these subsections is decision-eligible. Unless stated otherwise, all
development values are five-seed ensemble means on the 249,072 common keys, with
station-median RMSE in °C. Section 4.6 contains the evaluation-period slots.

\subsection{The reference model, not the architecture, sets the reported gain}\label{the-reference-model-not-the-architecture-sets-the-reported-gain}

How much of a reported gain survives a strong reference? Most of it does not. On
identical keys, the same fixed model reports a seven-day median station skill of
+0.251 against persistence and +0.038 against damped persistence --- a factor of
6.6 between two numbers that differ only in the denominator of equation (10).

{\small
\setlength{\tabcolsep}{3pt}
\renewcommand{\arraystretch}{1.12}
\noindent
\begin{tabularx}{\dimexpr\linewidth-24pt\relax}{@{}rrrrr>{\raggedright\arraybackslash}X@{}}
\toprule\noalign{}
Lead & Persistence & Damped persistence & LightGBM & LSTM & ThermoRoute \\
\midrule\noalign{}
1 d & 0.803 & 0.774 & 0.578 & 0.662 & 0.631 \\
3 d & 1.576 & 1.406 & 1.280 & 1.323 & 1.291 \\
7 d & 2.217 & 1.738 & 1.649 & 1.679 & 1.657 \\
\bottomrule\noalign{}
\end{tabularx}
}

\emph{Station-median RMSE, °C.} Read down the seven-day row: station-median RMSE falls
from 2.217 °C under persistence to 1.738 °C under damped persistence to 1.657 °C
under ThermoRoute. Of the 0.560 °C total reduction, 0.479 °C is delivered by
relaxing yesterday\textquotesingle s observation toward a fitted seasonal climatology, and
0.081 °C by everything the deep model adds. The same arithmetic in skill units
gives +0.203, +0.187, and +0.251 against persistence and +0.168, +0.076, and
+0.038 against damped persistence at 1, 3, and 7 days, all dimensionless.

The gap widens with lead, which is the diagnostic signature of damping rather
than of learned river behaviour: the longer the lead, the more of the error a
seasonal relaxation removes on its own, and the less remains for a learned model
to remove. A study reporting only the persistence column would describe this
model as gaining a quarter of the seven-day error budget. A study reporting both
columns describes the same fitted weights as gaining four percent.

Reported against the strong reference, the paired effects of equation (9) are
−0.130, −0.104, and −0.062 °C at 1, 3, and 7 days, with ThermoRoute favoured at
0.86, 0.90, and 0.93 of the 120 stations, and exploratory whole-HUC2
cluster-bootstrap intervals of {[}−0.177, −0.081{]}, {[}−0.126, −0.083{]}, and
{[}−0.071, −0.052{]} °C. Those intervals carry the 15-cluster structure discussed in
Section 3.6, and every row of the development mirror is flagged
\texttt{NO\_\allowbreak{}STRONG\_\allowbreak{}INFERENCE} with the warnings \texttt{SMALL\_\allowbreak{}CLUSTER\_\allowbreak{}COUNT\_\allowbreak{}LT\_\allowbreak{}30} and
\texttt{LOW\_\allowbreak{}EFFECTIVE\_\allowbreak{}CLUSTER\_\allowbreak{}FRACTION}. If the reference set determines the size of the
reported gain, the next question is whether it also determines which model wins.

\subsection{The most accurate model on this panel is a gradient-boosted tree}\label{the-most-accurate-model-on-this-panel-is-a-gradient-boosted-tree}

It does. The tree ensemble has the lowest development-period station-median RMSE
at all three leads: 0.578, 1.280, and 1.649 °C against ThermoRoute\textquotesingle s 0.631,
1.291, and 1.657 °C. The paired station-level differences
(\texttt{ThermoRoute\ \allowbreak{}−\ \allowbreak{}LightGBM}, equation 9) are +0.046 °C at 1 day, +0.008 °C at 3
days, and −0.005 °C at 7 days, with ThermoRoute win rates of 0.00, 0.40, and 0.60
across the 120 stations. The one-day win rate is 0.00: not one of the 120
stations favours the constrained architecture at the shortest lead.

We report this as the primary result it is. On this panel, on identical keys, a
well-tuned gradient-boosted tree with site identity is more accurate than the
constrained deep architecture, decisively at 1 day and negligibly at 3 and 7,
where the exploratory cluster-bootstrap intervals for the paired difference are
{[}+0.001, +0.020{]} and {[}−0.010, +0.010{]} °C. The global LSTM is intermediate at
0.662, 1.323, and 1.679 °C. The architecture\textquotesingle s remaining practical argument is
not accuracy but the algebraic bound of equation (7), which held on 100.00\% of
audited rows, with a maximum absolute correction of 1.0000 °C against the
configured 1 °C limit and the derived station-by-lead RMSE inequality holding in
all 360 cells --- a property of the construction rather than a fitted outcome.

\subsection{An information-matched plain convolutional network reproduces the architecture}\label{an-information-matched-plain-convolutional-network-reproduces-the-architecture}

If a tree is more accurate, is any of the architecture\textquotesingle s structure doing work?
Very little of it is. Against the information-matched plain causal temporal
convolutional network of Section 3.2 --- same inputs, same anchor, same optimiser
budget, 38,346 against 38,505 parameters --- the full architecture\textquotesingle s paired median
station effect, computed within seed on identical keys, ranges across the five
seeds from −0.0003 to −0.0105 °C at 1 day, −0.0077 to −0.0179 °C at 3 days, and
−0.0035 to −0.0226 °C at 7 days. Every seed favours the full model; no seed by
more than 0.023 °C. Against the plain multilayer perceptron the same effects
range from −0.0360 to −0.0411 °C at 1 day, so sequence structure matters and the
components layered on top of it do not.

The one-factor deletions agree. Removing the temporal encoder is the largest
single effect, moving 1-day station-median RMSE from 0.631 to 0.679 °C and 7-day
from 1.657 to 1.675 °C. Removing the router (0.634), the mixture (0.634), the
dynamic prior (0.628), the fixed-relaxation constraint (0.627), or the residual
bound (0.630) moves the 1-day figure by at most 0.004 °C in either direction.
Replacing everything by the damped anchor alone gives 0.774, 1.406, and 1.738 °C.
These are five-seed deletion and intervention sensitivities on paired keys; they
do not prove component necessity or a capacity-matched attribution. Read
together with Section 4.2, they locate this model\textquotesingle s accuracy in the causal
sequence encoder and in the anchor it corrects, not in the routing, mixture, or
bounding machinery.

\emph{Comparability note.} The controls in this subsection are per-seed values from
the development-controls run and are compared only with each other and with the
matched full-architecture arm from the same run (0.6452, 1.3047, 1.6682 °C).
They are not the five-seed ensemble means of Sections 4.1--4.2 and must not be
differenced against them.

\subsection{Whole-region holdout removes a third of the transfer skill a random split reports}\label{whole-region-holdout-removes-a-third-of-the-transfer-skill-a-random-split-reports}

How much of a reported spatial transfer survives when the map is partitioned
honestly? Between a quarter and a third of it does not.

{\small
\setlength{\tabcolsep}{3pt}
\renewcommand{\arraystretch}{1.12}
\noindent
\begin{tabularx}{\dimexpr\linewidth-24pt\relax}{@{}>{\raggedright\arraybackslash}X>{\raggedright\arraybackslash}Xrr>{\raggedright\arraybackslash}X@{}}
\toprule\noalign{}
Arm & Design & 1 d & 3 d & 7 d \\
\midrule\noalign{}
Temporal development & seen stations, 2019--2020 & +0.203 & +0.187 & +0.251 \\
Random held-site warm start & 4 folds; held-site history in preprocessing & +0.178 & +0.172 & +0.241 \\
Held-region gauged transfer & leave-HUC2-region-out & +0.155 & +0.116 & +0.155 \\
\bottomrule\noalign{}
\end{tabularx}
}

\emph{Median station skill versus persistence, dimensionless.} Against damped
persistence the same three arms give +0.168 / +0.076 / +0.038 (temporal),
+0.145 / +0.061 / +0.030 (random held-site), and +0.147 / +0.086 / +0.075
(held-region). The three-day skill against persistence falls from +0.187 to
+0.172 when sites are held out at random and to +0.116 when whole regions are,
and the seven-day figure falls from +0.251 to +0.241 to +0.155. The gap between
the second and third rows is the direct measure of how much a random spatial
split flatters a model on this panel: the random split retains 92\% of the
temporal-arm three-day skill, the regional split 62\%.

The mechanism is neighbourhood, and the geometry supports that reading. Under
whole-region holdout the mean distance from a held-out station to its nearest
training gauge is 289 km, whereas in the intact cohort 19 of 120 stations have
another retained station within 10 km and 38 share a hydrologic unit code with
another retained station. A random held-site split on such a panel measures
something much closer to interpolation between instrumented neighbours than to
transfer to an uninstrumented region.

The ranking of models is unchanged by the harder partition, and the margins
widen. In the held-region arm the global tree ensemble again attains lower
station-median RMSE than ThermoRoute (0.652 against 0.676 °C at 1 day, 1.391
against 1.428 at 3 days, 1.786 against 1.860 at 7 days), with paired median
differences of +0.031 {[}+0.023, +0.040{]}, +0.038 {[}+0.029, +0.050{]}, and +0.067
{[}+0.049, +0.080{]} °C and ThermoRoute win rates near 0.16. The global LSTM, with
its station embedding disabled in this arm, gives 0.679, 1.445, and 1.876 °C.
Since the spatial partition moves the answer this much, the next question is
whether the remaining skill is spatially and seasonally uniform.

\subsection{Skill is regionally uniform, and interval coverage is bought with width}\label{skill-is-regionally-uniform-and-interval-coverage-is-bought-with-width}

It is uniform, and this is the one place where the constrained model\textquotesingle s behaviour
is unremarkable in a useful way. Region-weighted skill against persistence --- the
mean of the 15 per-HUC2 medians --- is +0.204, +0.189, and +0.253 at 1, 3, and 7
days, essentially unchanged from the pooled medians of +0.203, +0.187, and
+0.251, so no single region carries the headline. Per-region medians against
persistence at 1 day range from +0.131 (HUC2:10) to +0.285 (HUC2:18) and at 7
days from +0.217 (HUC2:06) to +0.290 (HUC2:09). Stratifying by drainage area into
three groups of roughly 39 stations gives 1-day medians of +0.190, +0.214, and
+0.231 from small to large. Against damped persistence the seven-day per-region
medians compress to a range of +0.012 to +0.077 --- the same message as Section
4.1, restated spatially: what varies across regions is small once the seasonal
reference is doing its work.

Interval behaviour tells the complementary story about what calibration costs.
Split-conformal intervals on the 249,072 development keys attain an empirical
marginal coverage of 0.909 against a nominal 0.90, with a mean width of 3.87 °C
and a mean interval score of 4.93; per lead, coverage is 0.905, 0.910, and 0.912
at widths of 2.01, 4.22, and 5.38 °C. On the warm-season tail slice defined by
training-period q90 exceedance (30,531 keys), coverage is 0.918 at a width of
3.50 °C. The two pre-specified sensitivities show how much of that is a
convention. Calibrating on the maximum nonconformity score within blocks of seven
consecutive retained calibration rows raises overall coverage to 0.981 but widens
the interval to 5.74 °C --- coverage bought with width, not with sharpness. The
idealized delayed adaptive-conformal variants track nominal coverage more closely
as the step size grows (0.903, 0.900, and 0.892 for γ = 0.005, 0.02, 0.05) but
produce unbounded widths in several slices, so their interval scores are not
finite. We report width and interval score beside every coverage figure precisely
so that adaptive calibration is not presented as free, and none of these figures
is a conditional-coverage statement.

Input perturbations locate the model\textquotesingle s dependence in the same two channels the
ablations pointed to. Gaussian sensor noise at 0.25 and 0.5 training standard
deviations degrades 1-day RMSE from 0.631 to 1.621 and 2.999 °C (+147\% and
+360\%); missing-forcing blocks of 3, 7, and 14 days cost about +0.13 °C at 1 day
and +0.04 °C at 7 days; air-temperature offsets of ±2 training standard
deviations cost +0.11 to +0.15 °C at 1 day; and multiplying discharge by 0.5 or 2
changes 1-day RMSE by at most +0.003 °C. On this panel the predictor is dominated
by the water-temperature and air-temperature channels and is nearly insensitive
to discharge. These are synthetic data-corruption probes, not climate
projections, physically coherent scenarios, or deployment-safety tests.

\subsection{Evaluation-period results}\label{evaluation-period-results}

The following slots are rendered deterministically from the verified acquisition
receipt after the single opening described in Section 3.7. Each is reported
exactly once, including adverse, conflicting, and non-reportable outcomes, and
each carries the fixed-cohort descriptive verdict of Section 3.6. Handwritten
substitution into these slots is rejected by body hashes and evidence bindings.

\textbf{Table 4.1 --- the five formal rows.} Fields: model pair, lead, predeclared
margin, reportable station count, reportable cluster count, median paired station
RMSE difference, whole-HUC2 cluster-bootstrap 95\% interval, win rate, exact
sign-flip raw p-value, Holm-adjusted p-value, and the bound gate verdict.

{\small
\setlength{\tabcolsep}{3pt}
\renewcommand{\arraystretch}{1.12}
\noindent
\begin{tabularx}{\dimexpr\linewidth-24pt\relax}{@{}l>{\raggedright\arraybackslash}Xrl>{\raggedright\arraybackslash}X>{\raggedright\arraybackslash}X>{\raggedright\arraybackslash}X@{}}
\toprule\noalign{}
\# & Comparison & Lead & Margin & Effect and interval & Sensitivities (p\_raw / p\_Holm) & Verdict \\
\midrule\noalign{}
1 & ThermoRoute vs. damped persistence & 1 d & 0.00 °C & \texttt{{[}TO\ \allowbreak{}BE\ \allowbreak{}FILLED\ \allowbreak{}AFTER\ \allowbreak{}OPENING{]}} & \texttt{{[}TO\ \allowbreak{}BE\ \allowbreak{}FILLED\ \allowbreak{}AFTER\ \allowbreak{}OPENING{]}} & \texttt{DESCRIPTIVE\_\allowbreak{}ONLY\_\allowbreak{}INFERENCE\_\allowbreak{}GATE\_\allowbreak{}FAILED} \\
2 & ThermoRoute vs. damped persistence & 3 d & 0.00 °C & \texttt{{[}TO\ \allowbreak{}BE\ \allowbreak{}FILLED\ \allowbreak{}AFTER\ \allowbreak{}OPENING{]}} & \texttt{{[}TO\ \allowbreak{}BE\ \allowbreak{}FILLED\ \allowbreak{}AFTER\ \allowbreak{}OPENING{]}} & \texttt{DESCRIPTIVE\_\allowbreak{}ONLY\_\allowbreak{}INFERENCE\_\allowbreak{}GATE\_\allowbreak{}FAILED} \\
3 & ThermoRoute vs. damped persistence & 7 d & 0.00 °C & \texttt{{[}TO\ \allowbreak{}BE\ \allowbreak{}FILLED\ \allowbreak{}AFTER\ \allowbreak{}OPENING{]}} & \texttt{{[}TO\ \allowbreak{}BE\ \allowbreak{}FILLED\ \allowbreak{}AFTER\ \allowbreak{}OPENING{]}} & \texttt{DESCRIPTIVE\_\allowbreak{}ONLY\_\allowbreak{}INFERENCE\_\allowbreak{}GATE\_\allowbreak{}FAILED} \\
4 & ThermoRoute vs. LightGBM & 3 d & +0.05 °C & \texttt{{[}TO\ \allowbreak{}BE\ \allowbreak{}FILLED\ \allowbreak{}AFTER\ \allowbreak{}OPENING{]}} & \texttt{{[}TO\ \allowbreak{}BE\ \allowbreak{}FILLED\ \allowbreak{}AFTER\ \allowbreak{}OPENING{]}} & \texttt{DESCRIPTIVE\_\allowbreak{}ONLY\_\allowbreak{}INFERENCE\_\allowbreak{}GATE\_\allowbreak{}FAILED} \\
5 & ThermoRoute vs. LightGBM & 7 d & +0.05 °C & \texttt{{[}TO\ \allowbreak{}BE\ \allowbreak{}FILLED\ \allowbreak{}AFTER\ \allowbreak{}OPENING{]}} & \texttt{{[}TO\ \allowbreak{}BE\ \allowbreak{}FILLED\ \allowbreak{}AFTER\ \allowbreak{}OPENING{]}} & \texttt{DESCRIPTIVE\_\allowbreak{}ONLY\_\allowbreak{}INFERENCE\_\allowbreak{}GATE\_\allowbreak{}FAILED} \\
\bottomrule\noalign{}
\end{tabularx}
}

\textbf{Table 4.2 --- all-model scores on the evaluation common keys.}
\texttt{{[}TO\ \allowbreak{}BE\ \allowbreak{}FILLED\ \allowbreak{}AFTER\ \allowbreak{}OPENING{]}} --- station-median RMSE for every primary model at
every lead on the exact common key registry, with the reportable station count.
The frozen model registry for this cohort is the six primary models together
with the seven one-factor architecture controls of Section 3.2, so the control
rows are rendered here on the same keys and the same denominators; they remain
exploratory deletion and intervention sensitivities and cannot promote or
replace a formal row.

\textbf{Table 4.3 --- temporal coverage and calendar balance.}
\texttt{{[}TO\ \allowbreak{}BE\ \allowbreak{}FILLED\ \allowbreak{}AFTER\ \allowbreak{}OPENING{]}} --- for each of the five rows, all eight predeclared
descriptive candidates (equal-weighted 12 year-by-season cells; leave-one-year
for 2021, 2022, 2023; leave-one-season for DJF, MAM, JJA, SON) and the
deterministic most-adverse value.

\textbf{Table 4.4 --- interval and probability behaviour on the evaluation keys.}
\texttt{{[}TO\ \allowbreak{}BE\ \allowbreak{}FILLED\ \allowbreak{}AFTER\ \allowbreak{}OPENING{]}} --- by model and lead: empirical marginal coverage
at the nominal 90\% level, mean interval width, the equal-weight three-quantile
pinball mean, Brier score and Brier skill against the frozen seasonal reference,
log loss, AUROC, AUPRC, expected calibration error, calibration slope and
intercept, and station-balanced reliability bins. This table is filled from the
evaluation-period probabilistic family computed by the trusted scorer during the
single opening; it is labelled exploratory, consistent with Section 3.6, and it
is a different object from the development-period suite of Section 6.3. Three
fields are not produced by the opening and are rendered as explicit status
tokens rather than as numbers or blanks: the interval score, the block-maximum
calibration sensitivity, and the delayed adaptive-conformal sensitivity.

\textbf{Table 4.5 --- outcome quality control.}
\texttt{{[}TO\ \allowbreak{}BE\ \allowbreak{}FILLED\ \allowbreak{}AFTER\ \allowbreak{}OPENING{]}} --- immutable counts by cohort, station, variable,
raw qualifier string, and value presence; multiple-finite-series conflict counts;
the outcome quality-control gate status; and the qualifier-restricted \texttt{\{A\}}
sensitivity for each formal row.

\textbf{Table 4.6 --- external arm on the evaluation keys.}
\texttt{{[}TO\ \allowbreak{}BE\ \allowbreak{}FILLED\ \allowbreak{}AFTER\ \allowbreak{}OPENING{]}} --- the 30-site site-identifier-disjoint arm,
labelled exploratory and history-dependent. The leave-HUC2-region-out arm has no
evaluation-period counterpart: the opening produces no held-region artifact, so
that fragment renders as an explicit status token and the whole-region holdout
evidence in this paper remains the development-period evidence of Section 4.4.

\begin{center}\rule{0.5\linewidth}{0.5pt}\end{center}

\section{Discussion}\label{discussion}

\subsection{Why these numbers are smaller than published gains}\label{why-these-numbers-are-smaller-than-published-gains}

Our reported gains are far smaller than those commonly published for daily river
temperature, and the difference is a difference of design rather than of model
quality. Three specific choices account for it, and each can be read off Section
4 directly.

The first is the reference. Where a study quotes skill against persistence or
climatology, the number it reports is comparable to our +0.203 to +0.251 column,
not to our +0.038 to +0.168 column. On a variable with this much day-to-day
memory the two differ by up to a factor of 6.6 for the same fitted weights, and
the discrepancy grows with lead, so studies at longer leads are the ones most
affected. We would expect much of the between-study dispersion noted by
\href{https://doi.org/10.5194/hess-29-2521-2025}{Corona and Hogue (2025)} to be
absorbed by re-scoring published models against a fitted damped anchor rather
than against persistence, and we suggest damped persistence as the minimum
reference for this variable.

The second is the spatial partition. A model evaluated on randomly held-out
gauges in a dense network is being asked to interpolate between instrumented
neighbours; the same model evaluated on whole held-out regions, 289 km from the
nearest training gauge, retains 62\% rather than 92\% of its three-day skill. Where
published transfer results use random site splits, they are not measuring the
quantity their framing implies, and the difference is roughly a third of the
reported skill on this panel. This is a scale effect rather than a modelling
disagreement: at the spacing of a national monitoring network, random site
holdout and regional holdout are different experiments.

The third is the key set and the preprocessing boundary. Scoring each model on
its own complete-case set and fitting scaling, imputation, climatology, event
thresholds, or interval offsets on windows that include the evaluated interval
both move error downward without leaving a visible trace. We cannot quantify how
much of the published literature is affected, and we do not assert that it is;
what we can say is that removing both channels here left the constrained deep
model behind a gradient-boosted tree at every lead.

We are more cautious about the mechanistic reading of Section 4.3. That an
information-matched plain causal convolutional network reproduces the full
architecture to within 0.023 °C, and that deleting the router, the mixture, the
dynamic prior, or the residual bound moves 1-day error by at most 0.004 °C,
suggests that the accuracy available to a point-scale statistical predictor on
this panel is largely exhausted by a causal sequence encoder over issue-time
history plus a seasonal anchor. It does not identify why. Capacity matching is
not search-budget matching, deletion is not attribution, and the controls are
exploratory throughout.

\subsection{What a benchmark of this geometry can carry}\label{what-a-benchmark-of-this-geometry-can-carry}

The evaluation design of this study is stronger than its inferential reach, and
the two should not be confused. Clustered inference over a gauge panel rests on
resampling or sign-flipping whole spatial units, and the accuracy of those
procedures depends on having enough units, of comparable size, that the
asymptotics they invoke are not a fiction. Fifteen HUC2 groups with an
inverse-Herfindahl effective count of 9.54 and a largest group holding 21.7\% of
stations is not enough --- and, as Section 3.6 sets out, the cluster count is the
least binding of three independent reasons why this study is descriptive. The
sampling design was never exchangeable across regions, the null-simulation
calibration was never built, and the threshold that has attracted the most
attention was written at a value no U.S. cohort can reach.

The usual practice in applied water-temperature machine learning is to report a
station-level or row-level confidence interval, or a paired test across sites,
without addressing whether the number of independent spatial units can carry the
procedure. Under that practice a cohort like ours would produce intervals and
p-values that look authoritative, and nothing in the paper would tell a reader
that the effective number of clusters is under ten. The difference between that
paper and this one is not the data and not the model; it is whether the limit is
disclosed.

The design lesson is transferable. If a study intends region-clustered inference
over U.S. gauges, the cluster structure must be part of the \emph{sampling design},
not a post-hoc grouping of whatever cohort availability produced, and the
partition and its power must be fixed before any outcome is accessible. Reaching
a genuinely large number of independent spatial units requires either a finer
partition with a defensible independence argument or a deliberately stratified
draw across many more regions; neither is achievable by relabelling an existing
cohort. And a threshold stated in units of a fixed national partition should be
checked against how many such units exist before it is sealed. A second instance
of the same discipline appears in Section 6.3, where a pre-registered data
contract rejected an output we would have preferred to report and we left the
contract alone. A pre-registration that never binds against its authors\textquotesingle{}
convenience has not been tested; these two cases are the evidence that this one
has.

\subsection{What transfers}\label{what-transfers}

This study is narrow by construction. It is a common-key, clustered,
pre-registered benchmark of point predictors at gauged sites, and it is not an
attempt to replace process-based thermal models, river-network graph models, or
differentiable hybrid formulations
(\href{https://doi.org/10.1137/1.9781611976700.69}{Jia et al., 2021};
\href{https://doi.org/10.1029/2023WR034420}{Rahmani et al., 2023};
\href{https://doi.org/10.3389/frwa.2023.1184992}{Zwart et al., 2023}), nor
architectures that encode geographic context across regions and scales
(\href{https://doi.org/10.1145/3748636.3762716}{Luo et al., 2025}), each of which
represents structure --- connectivity, upstream forcing, energy balance, spatial
context --- that a point-scale statistical predictor does not. Every arm here,
including the whole-region and site-identifier-disjoint arms, consumes the target
site\textquotesingle s observed water temperature through the issue date (Section 3.4), so the
paper is silent on prediction at locations with no thermal record.

What is portable is not the architecture but three controls that do not depend on
it: scoring every model on an identical key registry so that no model benefits
from selective prediction; fitting every preprocessing and calibration statistic
strictly backwards in time; and partitioning space by whole hydrologic regions
rather than by site. None requires unusual computational resources, and each
moved a reported number in this study by more than the difference between the
competing architectures. On this evidence, benchmark design is a larger source of
variation in reported river-temperature skill than model design is.

\begin{center}\rule{0.5\linewidth}{0.5pt}\end{center}

\section{Limitations}\label{limitations}

\subsection{Machine-verifiable scope statements}\label{machine-verifiable-scope-statements}

The following statements are part of the evidence contract and remain in the
manuscript regardless of the eventual numerical outcome.

Route A is history-dependent and uses target-site water-temperature observations through each issue date; it does not establish ungauged prediction.

Route A does not establish independent river-network separation, physical river-network routing, or hydraulic travel time.

Route A is a one-shot retrospective historical-information evaluation, not an operational replay with archived as-issued predictor vintages or future NWP.

Route A uses daily-mean statistical thresholds and a numerical non-inferiority margin; neither has ecological, biological, or regulatory meaning.

Route A provides no physical, deployment, regulatory, or distribution-free safety guarantee, and failure to reject is not equivalence.

Route-A architecture controls and predictor sensitivities are descriptive or exploratory and do not identify causal mechanisms.

Route A evaluates a fixed availability-enriched 120-site cohort and is not nationally representative of all U.S. rivers or all calendar days.

Route A does not establish conditional coverage, and its equal-weight three-quantile pinball summary is not CRPS.

\subsection{Cohort, measurement, and design limitations}\label{cohort-measurement-and-design-limitations}

The station sample is availability-enriched rather than randomly drawn, and the
coverage thresholds that produced it used the 2019--2020 interval, so that
interval is not independent of cohort construction. The original provider bytes
for the 2006--2020 panel are unavailable, so development reproduction begins from
the committed derived artifact and the 1,465-candidate discovery execution is
auditable but not replayable; the panel retains no NWIS qualifier columns or
method and sensor history for that period, so measurement discontinuities cannot
be reconstructed for the development record. Meteorology is represented at the
station coordinate rather than integrated over the contributing catchment, which
is least appropriate for the largest basins. Two stations contain 2,059 signed
negative discharge rows whose semantics are unresolved. Outcomes are daily means,
and the evaluation is conditioned on outcome observability. HUC2 is a coarse
administrative grouping and is not an independent river-network component. The
+0.05 °C margin is a numerical ceiling with no stakeholder-derived importance.
Covariates are retrospectively acquired latest-provider values rather than
as-issued vintages, and model-selection budgets were documented but not equalized
across model classes. Training wall-clock time and inference throughput were not
recorded for the frozen suite, and no latency, memory, energy, or multi-hardware
benchmark is reported. The evidence chronology is owner-controlled, with no
external timestamp or independent custodian. The air2stream-style hybrid
reference was never fitted, so the hybrid process family is absent from every
comparison in this paper.

\subsection{The development-period probabilistic suite is not reported}\label{the-development-period-probabilistic-suite-is-not-reported}

This section concerns the development-period probabilistic evaluation only. The
evaluation-period probability family --- coverage, width, three-quantile pinball,
Brier score and skill, log loss, discrimination, calibration slope and intercept,
and station-balanced reliability bins --- is computed by a separate code path
inside the single opening and is reported in Table 4.4. Nothing below extends to
it, and the two must not be read as one withheld result.

The pre-registered development-period probabilistic contract treats a zero-width
nominal prediction interval as a fatal condition and prohibits evaluation-time
repair. In the development panel, 135 of 26,993,675 member-level rows carrying
complete quantile heads --- 0.0005\%, or about one row in 200,000 --- have identical
5th, 50th, and 95th percentile predictions. All 135 come from the tree ensembles (123 from the global
model, 12 from the per-station variant); no other model is affected. We emphasise
what the artifact is \emph{not}: there is no strict quantile-ordering violation
anywhere in the panel, and the maximum monotonicity violation is exactly
0.000 °C. The rejected rows are correctly ordered but degenerate.

Their distribution is interpretable, and we flag the reading that follows as
interpretation rather than as a measured attribution. The 135 rows fall on 12
distinct sites, all in hydrologic regions 04, 05, and 06 --- the Great Lakes, upper
Mississippi, and Missouri basins --- and 129 occur at the 1-day lead, 6 at 3 days,
and none at 7 days; in 95.6\% of them the observed target is below 1 °C. This is
the signature of ice-affected winter conditions, in which daily mean water
temperature is pinned near the freezing point and the conditional percentiles
genuinely coincide. The 12 per-station cases are constant-leaf predictions from a
model fitted on a small station-specific sample. The model outputs are
defensible; it is the contract\textquotesingle s inclusive equality test that refuses them.

We did not amend the contract. It was frozen before any access to the evaluation
labels, and relaxing it would have invalidated the frozen training receipts on
which the chronology gate of Section 3.7 depends. The development-period
probabilistic metric suite --- interval coverage from the nominal heads,
three-quantile pinball mean, reliability, and Brier score over 2019--2020 --- is
therefore not reported, and the analysis stage that consumes it is likewise not
reported. Two things follow. Development-period interval validity rests entirely
on the split-conformal analysis of Sections 3.5 and 4.5, which is unaffected: the
conformal offset is non-negative by construction and the deployed contract
separately requires positive width, so no delivered interval is degenerate. And
the standard event-probability verification metrics (Wilks, 2011) are unavailable
for the development period, so no development-period claim about probabilistic
calibration beyond empirical marginal interval coverage should be read into this
paper; the corresponding evaluation-period metrics are reported in Table 4.4.
Recovering the development-period suite requires a sealed erratum and a new
training lineage.

\subsection{Threshold and archive scope}\label{threshold-and-archive-scope}

An optional threshold analysis is separate from prediction evaluation. It can
only compare independently sourced observed daily maxima, expressed as a strict
seven-consecutive-day average of daily maxima, against a complete site-specific
standards registry, with seasonal applicability assigned by each window\textquotesingle s ending
date; season-external or incomplete windows remain unclassified, and ambiguous
standard matches fail closed. Any reported exceedance is descriptive: it is an
observation, not a model result and not a legal determination.

The current archive profile is machine-marked local-owner evidence: public
redistribution and third-party transfer are disabled and the public build mode
fails closed pending a byte-bound rights review, so this profile must not be read
as a public software or data release. The self-contained history bundle
deliberately retains reachable deleted objects so that chronology can be audited;
it is not a byte-level purge, and those objects are provenance rather than
current evidence.

\begin{center}\rule{0.5\linewidth}{0.5pt}\end{center}

\section{Conclusions}\label{conclusions}

We evaluated daily river water-temperature prediction at 120 stable U.S. gauges
across 34 states and 15 hydrologic regions over 2006--2020, scoring every model on
one common set of 249,072 station/\allowbreak date/\allowbreak horizon keys at 1-, 3-, and 7-day leads,
with every preprocessing and calibration statistic fitted strictly backwards in
time, and with the spatial partition varied from seen stations to random
held-site folds to whole held-out regions.

Three quantitative answers follow. Reported skill depends far more on the
reference model than on the architecture: median station skill at a seven-day
lead is +0.251 against persistence and +0.038 against damped persistence, and of
the 0.560 °C reduction in station-median RMSE from persistence to the deep model,
0.479 °C is delivered by seasonal damping alone. Model ranking does not favour
architectural constraint: a gradient-boosted tree with site identity has the
lowest station-median RMSE at all three leads and at every station at a one-day
lead, and an information-matched plain causal convolutional network reproduces
the full architecture to within 0.023 °C, while deleting the router, the mixture,
the dynamic prior, or the residual bound moves one-day error by at most 0.004 °C.
Reported spatial transfer depends on how the map is cut: three-day skill against
persistence falls from +0.187 to +0.172 under a random held-site split and to
+0.116 under whole-region holdout, at a mean nearest-training-gauge distance of
289 km, on a panel where 19 of 120 stations have a neighbour within 10 km.

The mechanism common to all three is that each design choice absorbs an
alternative explanation into the reported number: seasonal damping into a weak
reference, selective prediction into a model-specific key set, and spatial
interpolation into a random site split. For water-resource practice this matters
because reported skill is what a manager compares when choosing a tool. On this
evidence, a published river-temperature skill score is not interpretable without
its reference model and its spatial partition, and the incremental value of
architectural complexity over a strong statistical anchor and a well-tuned tree
is small enough that it should be re-measured under these controls before it is
relied upon.

The formal comparisons of this study are permanently descriptive. At most 15 HUC2
groups with an inverse-Herfindahl effective cluster count of about 9.54 cannot
pass the outcome-free inference gate, and two further components of that gate ---
the structural sampling assumptions and an unimplemented null simulation --- fail
independently of cohort size, so the five comparisons are fixed-cohort
descriptive effects under \texttt{DESCRIPTIVE\_\allowbreak{}ONLY\_\allowbreak{}INFERENCE\_\allowbreak{}GATE\_\allowbreak{}FAILED}, with
whole-HUC2 bootstrap intervals, exact sign-flip p-values, and Holm adjustments
retained only as assumption-conditional sensitivities. Exploratory and secondary
analyses remain reportable as such and are not erased by the gate. No result here
is a superiority, non-inferiority, equivalence, or parity statement, and none
generalises to a national or U.S.-river population.

\begin{center}\rule{0.5\linewidth}{0.5pt}\end{center}

\section*{Open Research Statement}

\textbf{Data availability.} The derived daily panel (\texttt{panel\_\allowbreak{}usgs\_\allowbreak{}120v2.parquet}), the stable station registry
(\texttt{station\_\allowbreak{}registry\_\allowbreak{}v1.csv}), the hydrologic-unit metadata snapshot, the
candidate-rejection ledger, and the frozen panel manifest that binds them are
deposited as one versioned dataset at \texttt{{[}DATA\ \allowbreak{}DOI\ \allowbreak{}TO\ \allowbreak{}BE\ \allowbreak{}MINTED{]}} under
\texttt{{[}DATA\ \allowbreak{}LICENCE\ \allowbreak{}TO\ \allowbreak{}BE\ \allowbreak{}ASSIGNED{]}}. The evaluation-period acquisition record --- exact
request and response bytes, series identifiers, approval qualifiers, final URLs,
retrieval timestamps, and content hashes --- is added to that deposit as a new
version after the acquisition described in Section 3.7.

Neither the DOI nor the data licence can be assigned before the byte-level rights
review described in Section 6.4 completes. The derived panel encodes values from
three providers whose terms differ, and a derived product does not inherit the
most permissive of its upstream terms. Until every object proposed for the
deposit carries a recorded, evidence-backed redistribution decision, the
project\textquotesingle s release tooling refuses to build a public archive. The deposit is
therefore described here as planned and specified, not as existing.

\textbf{Primary observational sources.} All observations are obtained from public providers and none is the property of
the authors. Daily-value water temperature and discharge come from the U.S.
Geological Survey National Water Information System
(\href{https://doi.org/10.5066/F7P55KJN}{U.S. Geological Survey, 2024}; parameter
codes 00010, 00060, and 00065 with statistic code 00003). Air temperature,
precipitation, the relative-humidity proxy, and daylight-period mean incoming
shortwave radiation come from Daymet V4 R1 at ORNL DAAC
(\href{https://doi.org/10.3334/ORNLDAAC/2129}{Thornton et al., 2022}). Wind speed
comes from gridMET (\href{https://doi.org/10.1002/joc.3413}{Abatzoglou, 2013}). Each
retains its provider\textquotesingle s own terms, citation requirement, and access route
independently of this manuscript.

\textbf{Material that cannot currently be redistributed.} For any provider object whose redistribution terms are unresolved at deposit
time, the deposit carries, in place of the bytes: the product identifier and
version, the exact request specification, the retrieval code, and a SHA-256
manifest of the bytes as retrieved, so a third party can re-acquire identical
inputs and verify them without relying on redistribution. The classes currently
defaulting to exclusion, and the two excluded outright --- the third-party AGU
LaTeX class, and three legacy CSV files retained in the repository archive from
an earlier three-station case study whose source and collection terms are
unrecorded --- are enumerated in SI16; the legacy material is not evidence for any
statement in this manuscript and its status is recorded in
\texttt{protocols/\allowbreak{}legacy\_\allowbreak{}three\_\allowbreak{}site\_\allowbreak{}semantics\_\allowbreak{}notice\_\allowbreak{}v1.md}. Original provider responses
for the 2006--2020 development panel were not retained, so development
reproduction begins from the committed derived artifact; that limitation does not
apply to the evaluation-period acquisition, for which exact bytes are archived.

\textbf{Software availability.} The analysis software, the frozen protocol and its sealed amendments, the
probability-metric erratum, the model-suite registry, the completion receipts,
the per-stage manifests, the environment probe, the verification entrypoints, the
deterministic result renderer that generates Section 4.6 from the acquisition
receipt, and the fully transitive Python 3.12 dependency lock with package hashes
are archived at \texttt{{[}SOFTWARE\ \allowbreak{}DOI\ \allowbreak{}TO\ \allowbreak{}BE\ \allowbreak{}MINTED{]}}, corresponding to release
\texttt{{[}RELEASE\ \allowbreak{}TAG\ \allowbreak{}TO\ \allowbreak{}BE\ \allowbreak{}ASSIGNED{]}} of the source repository at
\texttt{{[}REPOSITORY\ \allowbreak{}URL\ \allowbreak{}TO\ \allowbreak{}BE\ \allowbreak{}CONFIRMED{]}}. The project\textquotesingle s own source code is released
under the MIT licence (SPDX identifier \texttt{MIT}). That licence covers the source
code only: it does not license the observational data, the archived provider
responses, the third-party typesetting class, or any redistributed dependency
binary, each of which retains its own terms and is enumerated with those terms in
an archive notice file.

Analyses were run on Python 3.12 with NumPy, pandas, pyarrow, SciPy
(\href{https://doi.org/10.1038/s41592-019-0686-2}{Virtanen et al., 2020}),
scikit-learn (Pedregosa et al., 2011), LightGBM
(\href{https://proceedings.neurips.cc/paper/2017/hash/6449f44a102fde848669bdd9eb6b76fa-Abstract.html}{Ke et al., 2017}),
PyTorch (Paszke et al., 2019), and the USGS \texttt{dataretrieval} client
(\href{https://doi.org/10.5066/P94I5TX3}{Hodson and Hariharan, 2023}).

\textbf{Reproduction.} Sections 4.1--4.5 are reproducible from the archived panel, the archived model
bundles, and the pinned environment. Section 4.6 is reproducible from the
archived acquisition record and the same bundles. Bit-level equality is expected
for the tree models and agreement to a documented tolerance for the neural
members; the tolerances, the verification commands, and the expected digests are
listed in the Supporting Information. Reproduction has not yet been carried out
by an operator independent of the authors on an independent host, and this
statement is not a claim that it has.

\textbf{Placeholder inventory.} Seven bracketed identifiers appear above: \texttt{{[}DATA\ \allowbreak{}DOI\ \allowbreak{}TO\ \allowbreak{}BE\ \allowbreak{}MINTED{]}},
\texttt{{[}DATA\ \allowbreak{}LICENCE\ \allowbreak{}TO\ \allowbreak{}BE\ \allowbreak{}ASSIGNED{]}}, \texttt{{[}SOFTWARE\ \allowbreak{}DOI\ \allowbreak{}TO\ \allowbreak{}BE\ \allowbreak{}MINTED{]}},
\texttt{{[}RELEASE\ \allowbreak{}TAG\ \allowbreak{}TO\ \allowbreak{}BE\ \allowbreak{}ASSIGNED{]}}, \texttt{{[}REPOSITORY\ \allowbreak{}URL\ \allowbreak{}TO\ \allowbreak{}BE\ \allowbreak{}CONFIRMED{]}},
\texttt{{[}TRAINING\ \allowbreak{}WALL-CLOCK\ \allowbreak{}TO\ \allowbreak{}BE\ \allowbreak{}RECORDED{]}}, and \texttt{{[}INFERENCE\ \allowbreak{}COST\ \allowbreak{}TO\ \allowbreak{}BE\ \allowbreak{}RECORDED{]}}.
None is a real identifier or value, and none may be replaced by a reserved,
draft, example, or estimated figure. Together with the author block, the
corresponding-author details, and the Acknowledgments below, they are the
complete set of fields that must be closed before submission; each is tracked,
with the input required to close it, in \texttt{docs/\allowbreak{}WRR\_\allowbreak{}SUBMISSION\_\allowbreak{}CHECKLIST.md}.

\acknowledgments

{[}FUNDING TO BE COMPLETED --- each funder with its award number, or an explicit
statement that the work received no external funding.{]}

{[}COMPUTATIONAL RESOURCES TO BE COMPLETED --- the facility or facilities on which
the model suite was fitted.{]}

We acknowledge the U.S. Geological Survey for the National Water Information
System, the ORNL Distributed Active Archive Center for Daymet V4 R1, and the
Northwest Knowledge Network for gridMET. Acknowledgment of a data provider is not
an endorsement of this analysis by that provider.

\textbf{Competing interests.} {[}COMPETING INTERESTS TO BE COMPLETED --- a declaration is
required from every author, including the explicit statement that none exists if
that is the case.{]}

\textbf{Author contributions.} {[}CREDIT ROLES TO BE COMPLETED --- assign each author to
the applicable CRediT roles: conceptualization, methodology, software,
validation, formal analysis, investigation, resources, data curation, writing ---
original draft, writing --- review and editing, visualization, supervision, project
administration, funding acquisition. The signed intake form is
\texttt{docs/\allowbreak{}FAIR\_\allowbreak{}SUBMISSION\_\allowbreak{}READINESS\_\allowbreak{}AND\_\allowbreak{}TEMPLATES.md} §2.{]}

\section*{Supporting Information}

Supporting Information accompanies this manuscript and contains: the cohort and
registry description with the candidate-rejection ledger (SI01); the issue-time
information boundary, the product-compatibility bridge, and the isolation,
publication, durability, and adversary-model details of the replay and
acquisition machinery referred to in Sections 3.3 and 3.7 (SI02); model equations
and unit conventions (SI03); the sealed protocol, its amendments, and the
probability-metric erratum (SI04); the registered five-row comparison set and its
fields (SI05--SI06); all-model scores on the exact common keys, including the
status, provenance, and unofficial character of the unfitted air2stream-style
hybrid reference (SI07); probability and reliability schemas with the
non-reporting disposition of Section 6.3 and the full measured degenerate-interval
table (SI08); architecture and information-matched controls with seed and budget
provenance (SI09); the temporal coverage audit (SI10); spatial and leave-cluster
sensitivities, including the recomputed cluster geometry at HUC2, HUC4, HUC6, and
HUC8 (SI11); outcome quality control and qualifier evidence (SI12); the
history-dependent external arm (SI13); missingness and failure cases (SI14);
reproduction hashes, commands, and environment parity fields (SI15); and the
rights and data dictionary (SI16).

Ten supporting figures accompany these sections: Figures S1--S3 describe the
frozen cohort geometry, the temporal roles and issue-time information boundary,
and the model and calibration dataflow, and contain no evaluation-period
quantity; Figures S4--S8 and S10 expand Section 4.6 and are rendered only after
the acquisition of Section 3.7; and Figure S9 is an optional development-period
conformal-calibration sensitivity, labelled as such in the figure itself, which
must not be compared numerically with any evaluation-period figure.

Each figure carries evidence from exactly one period. Figures 1 and S1--S3 are
pre-opening structural material, Figures 2 and 4 and Figures S4--S8 and S10 are
evaluation-period, and Figure 3 and Figure S9 are development-period and say so
inside the figure. No value from one period is compared numerically with a value
from another.

% BIBLIOGRAPHY.  agujournal2025.cls loads apacite and sets
% \bibliographystyle{apacite} itself, so no \bibliographystyle is emitted here
% and no separate .bst has to be vendored.
%
% TODO(AUTHORS, docs/WRR_SUBMISSION_CHECKLIST.md item 6.4): the canonical
% Markdown still carries its citations as author-year prose, so this TeX contains
% no \cite or \citeA key.  \nocite{*} is a deliberate, temporary bridge: it
% emits the full reference list from ../references.bib, which checklist items
% 6.1 and 6.2 verified to be reconciled one-to-one with the in-text citation
% list, so every printed entry is in fact cited in the prose.  It is not a
% substitute for the citation-convention conversion, and it must be deleted in
% the same change that introduces real \cite/\citeA keys.
\nocite{*}
\bibliography{../references}

\end{document}
